\section{Compiler and Domain-Specific Language Foundations}\label{sec:compiler-dsl-foundations}

A compiler is a specialized software system that translates source code written in a high-level programming language
into a lower-level representation, such as machine code, bytecode, or, in the context of this thesis, a binary MIDI
file.
The compilation process is traditionally divided into several distinct phases to manage complexity and support
systematic validation~\cite{Aho_Lam_Sethi_Ullman_2006,Cooper_Torczon_2011}.
We apply that phased structure when describing Motivo; here we summarize the theory common to
general-purpose and domain-specific compilers.

\subsection{Lexical and Syntactic Analysis}\label{detail-subsec:lexical-and-syntactic-analysis}

The first phase, \textit{lexical analysis} (scanning), reads the raw character stream of the source file and groups
characters into meaningful sequences called \textit{tokens}.
For example, the string \texttt{tempo} might be converted into a keyword token, while \texttt{120} becomes an integer
literal token.
Regular expressions and finite automata provide the theoretical basis for recognizing token classes; in practice,
scanner generators automate this phase from a token specification~\cite{Aho_Lam_Sethi_Ullman_2006}.

The second phase, \textit{syntactic analysis} (parsing), takes the stream of tokens and verifies that they form a
valid sequence according to the formal grammar of the language.
Context-free grammars and pushdown automata model this step~\cite{Aho_Lam_Sethi_Ullman_2006}.
The parser organizes the flat token stream into a hierarchical tree structure known as an \textit{Abstract Syntax Tree}
(AST).
The AST captures the logical operations of the code (e.g., nesting a \texttt{play} statement inside a \texttt{loop}
block) while discarding irrelevant textual details like whitespace and semicolons.

\subsection{Semantic Analysis}\label{detail-subsec:semantic-analysis}

While the parser ensures that a sentence is grammatically correct, the \textit{semantic analyzer} ensures that it makes
logical sense.
This phase traverses the AST to perform checks that cannot be captured by a context-free
grammar~\cite{Cooper_Torczon_2011}.
The semantic phase therefore centers on two responsibilities.

\topic{Type checking.} Ensuring that operations are applied to compatible data types (e.g., preventing the
addition of a boolean to a note).

\topic{Symbol resolution.} Maintaining a \textit{symbol table} to ensure that variables and functions are
declared before they are used, and that they are accessed within their valid scope.

\subsection{Intermediate Representation and Code
Generation}\label{detail-subsec:intermediate-representation-and-code-generation}

After validation, the compiler often translates the AST into an \textit{Intermediate Representation} (IR).
An IR is a lower-level, often flattened representation of the program that is easier to optimize and translate than the
hierarchical AST~\cite{Cooper_Torczon_2011}.
In standard compilers this might involve generating three-address code.
In the context of music generation, lowering to an IR involves unrolling loops and resolving relative timings into
absolute timestamps.

Finally, the \textit{code generation} phase translates the IR into the target output format.
This phase handles the intricate binary serialization details, such as endianness encoding and protocol-specific
constraints, yielding the final artifact.

\subsection{Why Domain-Specific Languages?}\label{subsec:dsl-rationale}

Domain-specific languages are useful when a problem has stable concepts that are awkward to express directly
in a general-purpose language. Van Deursen, Klint, and Visser describe DSLs as languages that trade
generality for concise notations and domain-level abstractions~\cite{van_Deursen_Klint_Visser_2000}. Fowler
makes a similar software-engineering argument: a DSL can be worthwhile when its notation lets domain ideas be
expressed more directly than the host implementation language~\cite{Fowler_2010}.
We apply that argument to symbolic composition: Motivo's implementation is written in C++23, but the user writes musical
structures rather than C++ vectors of binary MIDI messages.
The Motivo language and \texttt{motivoc} compiler instantiate the phases outlined above.
